\section{Context Window Extension Techniques}
\label{sec:extension}

The fundamental challenge facing transformer-based large language models in processing extended sequences stems from the quadratic computational complexity of self-attention mechanisms. Standard attention computation scales as $O(n^2)$ in both time and memory with respect to sequence length $n$, imposing severe practical limitations on the maximum context window that can be processed during both training and inference. When sequence length increases by a factor of ten, computational requirements increase by a factor of one hundred. This quadratic bottleneck has motivated extensive research into techniques that extend the effective context window of pre-trained language models beyond their original training limits, with contemporary approaches achieving context windows extending from the early 2-4K token limitations to 128K tokens in production systems and beyond 2M tokens in research demonstrations.

This section examines the major categories of context window extension techniques, beginning with modifications to rotary position embeddings, proceeding through alternative positional encoding schemes, examining architectural and system-level innovations, exploring memory-augmented approaches, and concluding with continual pre-training strategies that have become standard practice in modern long-context language model development.

\subsection{RoPE-Based Scaling Methods}

Rotary Position Embedding has emerged as the dominant positional encoding method in modern large language models, adopted in flagship architectures including LLaMA, Mistral, Gemma, and GPT-J. RoPE encodes absolute position information while enabling models to learn relative position dependencies by representing positions through rotation matrices operating in 2D planes across the embedding dimensions \citep{su2024roformer}. The mathematical elegance of RoPE, which applies position-dependent rotations to query and key vectors, enables the dot product between vectors to naturally encode relative positional information through the geometric properties of rotation. However, despite these theoretical advantages, RoPE-based models exhibit severe degradation when processing sequences longer than those encountered during training, a phenomenon termed the position extrapolation problem. The core challenge arises because RoPE's frequency components are optimized for the positional range seen during pre-training, and extending beyond this range introduces out-of-distribution rotational angles that cause attention score instabilities and catastrophic performance collapse.

\subsubsection{Position Interpolation}

\citet{chen2023extending} introduced Position Interpolation as a solution to the extrapolation problem by proposing a conceptually simple but theoretically grounded approach. Rather than extrapolating to unseen positional indices beyond the training range, Position Interpolation linearly down-scales input position indices to compress them into the original context window size. Formally, if a model was pre-trained on positions $[0, L]$ and must process a sequence of length $sL$ where $s > 1$ is the scaling factor, Position Interpolation applies the transformation $p' = p/s$ to map position $p$ to its interpolated value $p'$. This transformation ensures that the maximum position index matches the pre-training limit, keeping all positional encodings within the distribution seen during training.

The theoretical foundation for Position Interpolation rests on the observation that interpolation is fundamentally more stable than extrapolation in the frequency domain. \citet{chen2023extending} demonstrated that the theoretical upper bound on interpolation error is approximately 600 times smaller than the corresponding extrapolation error bound, providing mathematical justification for why interpolation enables reliable context extension with minimal fine-tuning. In practice, Position Interpolation enabled LLaMA models to extend from their original 2K token context to 32,768 tokens after fine-tuning for fewer than 1,000 training steps on long sequences. This remarkable efficiency stemmed from the fact that the model's attention patterns, learned during pre-training, remained valid under interpolation because no out-of-distribution positional encodings were introduced.

However, Position Interpolation suffers from a critical limitation that becomes apparent when examining its effect on RoPE's frequency components. By uniformly scaling all positions by the same factor $1/s$, Position Interpolation effectively compresses the entire positional space, which in the frequency domain corresponds to removing high-frequency components of the positional encoding. This uniform compression makes adjacent tokens appear perceptually closer in the model's positional space, potentially conflating tokens that should be distinguished. The degradation worsens as the scaling factor increases, with empirical results showing that Position Interpolation typically achieves reliable extensions of 2-4$\times$ the original context length before significant performance degradation occurs.

\subsubsection{NTK-Aware Scaling}

The limitations of Position Interpolation motivated the development of frequency-aware scaling methods that treat different frequency components of RoPE with varying degrees of interpolation. NTK-Aware scaling, proposed by community researcher bloc97 through GitHub and Reddit discussions, represents a breakthrough in context extension efficiency by recognizing that different frequency dimensions of RoPE embeddings serve different representational purposes and should be scaled accordingly. The method derives its name from the Neural Tangent Kernel perspective on how neural networks generalize, though the practical implementation focuses on dimension-dependent frequency scaling.

The core innovation of NTK-Aware scaling lies in its modification of the RoPE base frequency. While Position Interpolation scales positions uniformly, NTK-Aware scaling instead modifies the base frequency $\theta$ used in RoPE's sinusoidal functions according to $\theta_{new} = \theta_{original} \cdot \alpha^{d/(d-2)}$, where $\alpha$ is the desired extension factor and $d$ is the dimensionality. This modification distributes interpolation pressure across frequency dimensions in a non-uniform manner. High-frequency dimensions, which encode fine-grained local positional distinctions between nearby tokens, receive less interpolation pressure and thus preserve their ability to distinguish adjacent positions. Low-frequency dimensions, which encode coarse-grained global structure and long-range positional relationships, receive more interpolation pressure to accommodate the extended context.

This frequency-aware approach addresses Position Interpolation's fundamental weakness by preventing the loss of high-frequency information that is critical for local coherence. The practical impact is substantial: NTK-Aware scaling enables 4-8$\times$ context extension without any fine-tuning whatsoever, compared to Position Interpolation's maximum 2-4$\times$ extension even with fine-tuning. Perhaps most remarkably, NTK-Aware scaling requires modifying only three lines of code in the RoPE implementation, making it trivially easy to integrate into existing models. Empirical evaluations demonstrated that models extended with NTK-Aware scaling maintain substantially lower perplexity across extended context lengths compared to Position Interpolation, validating the importance of frequency-aware scaling strategies.

The community further refined NTK-Aware scaling through the NTK-by-parts extension, which applies different interpolation strategies to different parts of the frequency spectrum rather than using a single smooth function across all frequencies. This refinement addressed catastrophic failures observed with pure NTK-Aware scaling at certain specific context lengths, discovering that different frequency ranges benefit from distinct scaling approaches. The success of NTK-Aware scaling in enabling zero-shot context extension without fine-tuning demonstrated that the position extrapolation problem could be largely solved through careful frequency domain manipulation alone.

\subsubsection{YaRN}

Building upon the insights from NTK-Aware scaling, \citet{peng2024yarn} introduced YaRN (Yet another RoPE extensioN method), a compute-efficient approach that combined sophisticated frequency-aware scaling with attention temperature adjustments. YaRN represents the culmination of community-driven research on RoPE scaling, incorporating lessons learned from Position Interpolation, NTK-Aware scaling, and their variants into a unified framework that achieves state-of-the-art efficiency for context extension.

YaRN's fundamental innovation lies in its three-way classification of RoPE frequency dimensions into distinct groups that receive tailored treatment. High-frequency dimensions, which encode local positional patterns and fine-grained distinctions, are not interpolated at all, preserving the model's ability to maintain local coherence and distinguish between adjacent tokens. Low-frequency dimensions, which encode global structural information and long-range dependencies, receive full interpolation to accommodate the extended context range. Mid-range frequencies, which operate at intermediate scales, follow NTK-Aware scaling rules that smoothly transition between the high and low-frequency behaviors. This tripartite classification enables YaRN to preserve both local and global positional information while extending context windows.

In addition to frequency-aware scaling, YaRN introduces attention temperature scaling to address the phenomenon of attention entropy collapse during context extension. As context length increases, attention distributions can become overly concentrated or overly diffuse, degrading the model's ability to focus on relevant information. YaRN's temperature adjustment counteracts this effect by modulating attention sharpness as a function of context length, maintaining stable attention patterns across varying sequence lengths.

The computational efficiency of YaRN represents perhaps its most significant practical contribution. \citet{peng2024yarn} demonstrated that YaRN requires 10$\times$ fewer tokens and 2.5$\times$ fewer training steps compared to previous context extension methods. In concrete terms, YaRN extended LLaMA 2 models (both 7B and 13B parameter variants) to 128K token context windows after training for merely 400 steps with batch size 64 on the PG19 dataset. This corresponded to training on less than 0.1\% of the original pre-training data, yet achieved state-of-the-art performance on long-context benchmarks. Furthermore, YaRN demonstrated extrapolation capability beyond the fine-tuning dataset's context length, indicating that the method instilled robust long-context understanding rather than simply memorizing longer training sequences.

The practical impact of YaRN has been substantial, with the method being adopted in numerous production language models including Qwen, DeepSeek, and various open-source LLaMA derivatives. Its combination of effectiveness, efficiency, and relative simplicity of implementation has made it the de facto standard for RoPE-based context extension in the research community and industry applications.

\subsubsection{LongRoPE and Advanced Extensions}

\citet{ding2024longrope} pushed the boundaries of RoPE-based context extension to unprecedented lengths, demonstrating the feasibility of extending context windows to 2,048K tokens (over 2 million tokens) through three complementary innovations. LongRoPE's first contribution involves exploiting non-uniformities in positional interpolation requirements through evolutionary search. Rather than applying uniform interpolation rules across all positional dimensions and all regions of the extended context, LongRoPE identifies that different frequency components and different positional ranges benefit from different interpolation strategies. The method employs evolutionary search guided by perplexity on long-context validation data to discover optimal non-uniform interpolation schedules, providing superior initialization for subsequent fine-tuning. In scenarios without fine-tuning, this non-uniform interpolation alone enables 8$\times$ context extension.

LongRoPE's second innovation introduces a progressive extension strategy that stages the context expansion process into multiple phases. Rather than directly extending from the base context length to the target 2M token context in a single step, LongRoPE first fine-tunes the model to an intermediate length of 256K tokens. The model is then subjected to a second round of positional interpolation and fine-tuning to reach the final 2,048K token context window. This staged approach allows the model to gradually adapt to longer contexts, mitigating the distribution shift that occurs when abruptly extending context by orders of magnitude. The progressive strategy also enables more efficient exploration of the extended positional space during fine-tuning, as the model can leverage intermediate-length representations when learning to process even longer sequences.

The third key innovation of LongRoPE addresses a pervasive problem across context extension methods: degradation of performance on the original short context lengths. When models are extended to very long contexts through interpolation and fine-tuning, they often suffer reduced performance on the shorter sequences they were originally trained on. This phenomenon, sometimes termed "short-context forgetting," occurs because the fine-tuning process optimizes for long-context performance and may shift the model's representations in ways that harm short-context processing. LongRoPE introduces a short context recovery mechanism that readjusts the model on 8K length sequences after long-context extension, explicitly recovering performance at the original context scale. This recovery phase is relatively brief compared to the long-context fine-tuning but proves essential for maintaining balanced performance across all sequence lengths.

LongRoPE was successfully integrated into Microsoft's Phi-3 production model, demonstrating its practical viability beyond research demonstrations. The method achieves its impressive 2M token context window with remarkably modest fine-tuning requirements, requiring at most 1K fine-tuning steps at training lengths within 256K. Extensive experiments on LLaMA-2 and Mistral models across diverse tasks validated LongRoPE's effectiveness while confirming that it maintains compatibility with the original transformer architecture, requiring only minimal modifications to the positional embedding mechanism.

The evolution of RoPE scaling methods continued with recent work including LongRoPE2, which further refined the approach through evolutionary search guided by what the authors term "needle-driven" perplexity. This variant uses targeted perplexity measurements focused on specific retrieval tasks to guide the search for optimal interpolation parameters. LongRoPE2 also introduced mixed context window training, where models are trained on a mixture of sequences of varying lengths rather than exclusively on long sequences. This mixed training approach addresses short-context degradation more directly by ensuring continued exposure to short sequences throughout the long-context adaptation process. Additional innovations included fine-tuning model weights to adopt rescaled RoPE specifically for long sequences while preserving short-context performance, achieving what the authors characterized as near-lossless scaling across the full range of context lengths from short to ultra-long.

Beyond the mainstream RoPE scaling methods, recent research has explored more fundamental modifications to rotary embeddings. ComRoPE \citep{yu2025comrope}, published at CVPR 2025 despite being positioned primarily for vision tasks, generalized RoPE through trainable commuting angle matrices. This work provided theoretical analysis demonstrating that pairwise commutativity of rotation matrices is essential for positional robustness and scalability, offering new mathematical insights into why RoPE functions effectively and how its parameters can be learned rather than fixed. Such theoretical advances continue to deepen understanding of the mathematical foundations underlying practical context extension methods.

\subsection{Alternative Position Encoding Methods}

While RoPE-based scaling methods have dominated recent context extension research, alternative approaches to position encoding offer fundamentally different perspectives on the extrapolation problem. The most prominent alternative, ALiBi (Attention with Linear Biases), eschews positional embeddings entirely in favor of a recency-biased attention mechanism.

\subsubsection{ALiBi and Linear Bias Approaches}

\citet{press2022train} introduced ALiBi with the explicit goal of enabling the "train short, test long" paradigm, where models trained on relatively short sequences could extrapolate to substantially longer sequences at test time without fine-tuning. ALiBi eliminates positional embeddings entirely, neither adding position information to input embeddings at the model's input nor incorporating learned positional parameters anywhere in the architecture. Instead, ALiBi modifies the attention mechanism itself by applying a bias directly to the attention scores after the query-key dot product.

Formally, ALiBi modifies the attention computation as $\text{softmax}(\mathbf{q}_i \mathbf{K}^\top + m \cdot [-(i-1), \ldots, -2, -1, 0])$, where $m$ is a head-specific scalar slope that is fixed before training and remains constant during inference. The bias term creates a linear penalty proportional to the distance between query position $i$ and each key position, with the penalty increasing linearly for more distant tokens. Each attention head receives a different slope $m$, with slopes typically following a geometric progression (e.g., $\{1/2, 1/4, 1/8, 1/16, \ldots\}$) to enable different heads to specialize in different distance ranges.

The design of ALiBi encodes an inductive bias toward recency, explicitly penalizing attention to distant tokens regardless of their semantic relevance. This recency bias reflects an assumption that more recent tokens are generally more relevant for predicting or processing the current token, an assumption that holds reasonably well for many natural language processing tasks. The linear distance penalty creates a smooth degradation of attention weights as a function of distance, rather than requiring the model to learn distance-dependent attention patterns from the positional encodings.

The extrapolation capabilities of ALiBi proved remarkable in empirical evaluations. \citet{press2022train} demonstrated that a 1.3B parameter model trained on sequences of 1,024 tokens could extrapolate to 2,048-token sequences, achieving perplexity equivalent to sinusoidal position embedding models that had been trained directly on 2,048-token sequences. Beyond mere extrapolation to 2$\times$ the training length, ALiBi models maintained reasonable performance on sequences 5-10$\times$ longer than their training length, a degree of extrapolation that proved impossible for standard positional encodings. Furthermore, ALiBi provided practical training benefits, achieving 11\% faster training speed and 11\% reduced memory usage compared to models using learned positional embeddings, because the bias computation is simpler than embedding lookups and additions.

However, ALiBi's design involves tradeoffs that limit its applicability to certain scenarios. The recency bias, while beneficial for many language modeling tasks, may not be optimal for tasks requiring equal attention to distant context regardless of recency. Retrieval tasks, where relevant information may appear anywhere in a long context with equal probability, potentially suffer under ALiBi's distance penalty. In contrast, RoPE's learned rotational encodings do not impose an explicit recency bias, instead allowing the model to learn task-appropriate attention patterns. This distinction explains why RoPE has been preferred for general-purpose large language models despite ALiBi's superior zero-shot extrapolation capability.

Theoretical analysis comparing interpolation and extrapolation approaches has illuminated the fundamental tradeoffs involved. Interpolation methods like Position Interpolation and YaRN require selecting a target context length and a corresponding scaling factor during the adaptation process. This fixed scaling factor introduces a tradeoff: extending context capability at long lengths necessarily degrades performance at the original short lengths due to the compression of positional information. Extrapolation methods like ALiBi naturally handle arbitrary lengths without modification, but require learned parameters (the head-specific slopes) that must be appropriate for the range of lengths encountered during training and the range needed at inference. Dynamic scaling approaches, discussed earlier in the context of NTK-Aware scaling, attempt to bridge these approaches by adapting interpolation factors dynamically based on the actual sequence length at inference time.

\subsubsection{Advanced Positional Encoding Analysis}

Recent research has applied sophisticated analytical frameworks to understand positional encoding behavior during context extension. Distributional analysis of RoPE reveals that existing extension methods can introduce out-of-distribution rotary angles even when using interpolation, because the interpolation may be applied suboptimally across different frequency dimensions \citep{zhang2024distributional}. This distributional perspective motivated approaches that selectively apply interpolation or extrapolation on a per-dimension basis depending on measured disturbance scores that quantify how far each dimension's angles deviate from their training distribution. Such distributional minimization approaches demonstrated 4.33\% improvement over standard Position Interpolation when extending LLaMA2-7B to 16K tokens, illustrating the value of fine-grained per-dimension analysis.

Other work has explored decomposition of positional vectors into frequency-dependent components, enabling independent decisions about interpolation versus extrapolation for each frequency band. This component-wise analysis provides a principled framework for understanding why certain frequency bands benefit from interpolation while others prefer extrapolation or no modification at all. The resulting flexibility in positional encoding design has advanced the theoretical understanding of context extension while also yielding practical methods that outperform simpler uniform approaches.

\subsection{Architectural and System-Level Approaches}

While positional encoding modifications address context extension through representational changes, architectural and system-level innovations approach the problem through fundamental reconsideration of attention computation itself. These methods recognize that position encoding, while important, represents only one aspect of the context extension challenge, and that substantial gains can be achieved through memory-efficient attention computation, distributed processing, and hardware-aware optimization.

\subsubsection{Ring Attention and Distributed Context Processing}

\citet{liu2023ring} introduced Ring Attention as a fundamentally different approach to enabling ultra-long contexts, one that sidesteps position encoding limitations entirely by instead distributing attention computation across multiple devices. Ring Attention leverages blockwise computation of self-attention and feedforward layers, decomposing the attention mechanism to operate on sequence blocks rather than requiring the entire sequence to reside in a single device's memory. The devices are arranged in a conceptual ring topology, with each device holding a portion of the sequence and computing attention for its local blocks.

The key innovation of Ring Attention lies in its overlapping of communication and computation. During the inner loop of attention computation, each device computes attention between its local queries and its current key-value blocks. Simultaneously, it sends its key-value blocks to the next device in the ring topology and receives key-value blocks from the previous device. By carefully orchestrating the timing of computation and communication, Ring Attention ensures that computation on one block of key-value pairs proceeds in parallel with the transfer of the next block. As long as the computation time for processing a block exceeds the communication time for transferring a block, the overlapping results in zero added overhead compared to standard attention that processes the entire sequence on a single device.

The scalability implications of Ring Attention are profound. The approach enables sequences up to device-count times longer than can be processed with memory-efficient transformers on a single device, effectively multiplying the maximum context length by the number of available devices. \citet{liu2023ring} demonstrated training on sequences exceeding 100 million tokens without any approximations to the attention mechanism, representing sequences 500$\times$ longer than prior state-of-the-art memory-efficient methods. The approach maintains exact attention computation rather than resorting to sparse or approximate attention, preserving the full representational capacity of the transformer architecture.

Ring Attention found practical application in the Large World Model (LWM), which employed the technique for million-length vision-language training. The ability to process such extreme context lengths enabled new applications in domains requiring dense, long-range modeling, such as video understanding, long document processing, and extended dialogue modeling. The distributed nature of Ring Attention also provides natural parallelism for training, though it requires careful implementation to achieve the overlapping communication and computation that eliminates overhead.

Distributed extensions of memory-efficient attention continue to advance the field. DistFlashAttention builds upon FlashAttention (discussed below) by extending its memory-efficient computation to distributed settings. DistFlashAttention introduces token-level workload balancing to ensure even distribution of computation across devices, overlaps key-value communication with attention computation similar to Ring Attention, and incorporates rematerialization-aware gradient checkpointing to minimize memory usage during backpropagation. These innovations enabled DistFlashAttention to process sequences 8$\times$ longer than Ring Attention alone, achieving 4.45-5.64$\times$ speedup over Ring Attention on benchmark tasks. The method successfully trained LLaMA-7B with sequences up to 512K tokens, demonstrating the feasibility of distributed training at extreme context lengths.

\subsubsection{FlashAttention: Memory-Efficient Exact Attention}

While Ring Attention addressed context scaling through distribution across devices, FlashAttention tackled the problem through within-device memory optimization. \citet{dao2022flashattention} observed that the primary bottleneck for attention computation on modern GPUs is not the raw number of floating-point operations but rather the memory movement between different levels of the memory hierarchy. GPUs feature fast on-chip SRAM with limited capacity and slower high-bandwidth memory (HBM) with substantially larger capacity. Standard attention implementations materialize the full $N \times N$ attention matrix in HBM, requiring repeated transfers between HBM and SRAM that dominate the total computation time.

FlashAttention eliminates this memory bottleneck through an IO-aware algorithm that uses tiling to restructure attention computation. Rather than computing the full attention matrix and then applying softmax and value multiplication, FlashAttention divides the query, key, and value matrices into blocks that fit in SRAM, computes attention for each block, and incrementally combines the results. The key algorithmic innovation involves maintaining running statistics for the softmax normalization, enabling correct computation of attention outputs without ever materializing the full attention matrix. This restructuring reduces the number of memory reads and writes between HBM and SRAM from $O(N^2)$ to $O(N^2 d^2 M^{-1})$, where $d$ is the head dimension and $M$ is SRAM size, yielding substantial reductions for typical parameters.

The practical impact of FlashAttention on training and inference proved substantial. On GPT-2, FlashAttention achieved 7.6$\times$ speedup compared to standard PyTorch attention implementations. For BERT-large with 512-token sequences, FlashAttention provided 15\% speedup over the MLPerf 1.1 training record, which represented highly optimized baseline implementations. Beyond speed improvements, FlashAttention's memory efficiency enabled training with longer sequences and larger batch sizes. On GPT-2, FlashAttention enabled training sequences 4$\times$ longer than was possible with standard attention implementations for a given memory budget, translating to substantially reduced perplexity (0.7 improvement) due to the benefits of longer training contexts.

The evolution of FlashAttention through subsequent versions demonstrated continued optimization potential. FlashAttention-2 \citep{dao2023flashattention2} achieved an additional 2$\times$ speedup over the original FlashAttention through algorithmic improvements in parallelism and work partitioning, reaching 50-73\% of the theoretical maximum FLOPs per second achievable on A100 GPUs. This represented exceptional hardware utilization, approaching the limits of what is theoretically possible given memory bandwidth constraints. FlashAttention-3, optimized for Hopper-generation GPUs (H100), introduced asynchronous execution to better overlap different operations and added support for low-precision computation (FP8) to further accelerate processing while maintaining acceptable accuracy.

Hardware-specific optimizations have extended beyond FlashAttention's core algorithm. S2-Attention introduced heterogeneous context sharding across attention heads, where different heads attend to different subsets of the context tokens while collectively the full context is covered. This approach exploits the observation that attention heads in transformers often specialize in different types of patterns, and not every head needs access to every token. By sharding context across heads, S2-Attention reduced memory bandwidth requirements and achieved 8.79$\times$ to 25.3$\times$ wall-clock speedup compared to FlashAttention-2 while maintaining full attention performance. The method demonstrated perfect retrieval accuracy on 128K context benchmarks, confirming that strategic context sharding does not sacrifice representational capacity when designed appropriately.

Inference-time efficiency has motivated additional attention innovations. Recycled Attention addresses the memory burden of the key-value cache that must be maintained during autoregressive generation. As generation proceeds, the KV cache grows linearly with the number of generated tokens, and attention computation time grows quadratically. Recycled Attention alternates between full attention steps that attend to all past tokens and recycled attention steps that attend only to a reduced cache of important tokens. This alternating pattern balances computational efficiency with semantic preservation, maintaining high-quality generation while dramatically reducing the memory footprint and attention computation cost for long-context generation.

\subsection{Memory-Augmented and Streaming Methods}

An alternative paradigm for handling long contexts involves augmenting transformer models with explicit memory mechanisms that operate outside the standard attention framework. These memory-augmented approaches often draw inspiration from cognitive science and neural memory systems, introducing specialized components designed specifically for long-term information retention.

\subsubsection{StreamingLLM and Attention Sinks}

\citet{xiao2024efficient} discovered a surprising phenomenon in pre-trained large language models that enabled infinite-length sequence processing without fine-tuning. Their analysis revealed that language models exhibit what they termed "attention sinks," where initial tokens in a sequence receive disproportionately high attention weights regardless of their semantic content. This phenomenon occurs because the softmax attention mechanism requires weights to sum to one, and when there is no particularly relevant token to attend to, models learn to allocate attention to consistent sink tokens rather than distributing it diffusely across all tokens.

StreamingLLM exploits this attention sink phenomenon by retaining the initial tokens (typically 4-16 tokens) in the key-value cache alongside recent context maintained through a sliding window. By keeping these sink tokens permanently in the cache, StreamingLLM enables models to maintain stable attention distributions even as the context extends far beyond the training length. The model can attend to the sink tokens when needed to "dump" attention weights, preventing the attention distribution instabilities that would otherwise occur in infinite-length streaming scenarios.

The practical performance of StreamingLLM proved remarkable across multiple model families. StreamingLLM enabled LLaMA-2, MPT, Falcon, and Pythia models to process sequences exceeding 4 million tokens with stable performance, despite these models having been trained on contexts orders of magnitude shorter. Compared to naive sliding window approaches that recompute attention for the entire window at each step, StreamingLLM achieved 22.2$\times$ speedup in streaming settings by maintaining the KV cache for recent tokens and only computing attention for new tokens. This efficiency gain made real-time processing of effectively infinite-length streams computationally feasible.

The generality of the attention sink phenomenon across diverse model families and architectures suggests it represents a fundamental aspect of how softmax attention handles ambiguous attention scenarios. The discovery has influenced subsequent research on efficient attention mechanisms and memory-augmented models, highlighting the importance of understanding emergent behaviors in trained models rather than relying solely on architectural design principles.

\subsubsection{Landmark Attention and Block-Level Retrieval}

\citet{mohtashami2023landmark} introduced Landmark Attention as an approach to enable random-access to arbitrarily long contexts through learned retrieval of relevant blocks. The method divides input sequences into blocks and associates each block with a landmark token that serves as a compressed representation of the block's content. During attention computation, the model first attends to the landmark tokens to identify which blocks are relevant to the current query, then attends to the full content of the selected blocks.

This two-stage attention process enables efficient processing of very long contexts by avoiding the need to compute attention over every token in the full context. Instead, the model performs coarse-grained retrieval at the block level through landmark attention, followed by fine-grained attention within selected blocks. The landmark tokens are trained through the standard language modeling objective, learning to encode information that enables effective block selection. This training approach differs from retrieval-augmented methods that use separate retrieval systems, instead integrating retrieval directly into the attention mechanism.

Fine-tuning LLaMA-7B with Landmark Attention extended context to 32K tokens, matching the context capability of models like GPT-4 despite LLaMA's original 2K context limit. The method maintained accuracy on retrieval tasks while providing computational efficiency through block-level routing. The approach demonstrated that learned attention-based retrieval could provide an alternative to both full dense attention and external retrieval systems, with the routing decisions emerging from the model's training rather than being hard-coded.

\subsubsection{Recurrent Memory Transformers}

The Recurrent Memory Transformer (RMT) introduced by \citet{bulatov2022recurrent} combined transformer attention with recurrent connections across segments through special memory tokens. RMT processes long sequences by dividing them into segments, processing each segment with a transformer, and passing memory tokens from one segment to the next to maintain information flow. These memory tokens function as a learned recurrent state that aggregates information across segment boundaries, enabling the model to maintain long-term dependencies despite processing sequences in bounded-length chunks.

The key advantage of RMT lies in its linear scaling of computation with sequence length, achieved by maintaining constant segment length while increasing the number of segments for longer sequences. Memory tokens provide additional capacity beyond the segment length for cross-segment information aggregation, with the number and dimensionality of memory tokens serving as hyperparameters that control the information bottleneck between segments. Empirical evaluations demonstrated that RMT outperformed Transformer-XL on tasks requiring long-term dependencies, validating the effectiveness of learned memory tokens for sequential information propagation.

Extensions to the RMT framework have explored more sophisticated memory mechanisms. The Associative Recurrent Memory Transformer (ARMT) augmented RMT with associative memory mechanisms inspired by computational neuroscience, enabling more flexible information storage and retrieval patterns within the memory tokens. These associative mechanisms allow memories to be stored and retrieved based on content similarity rather than solely through sequential propagation, potentially enabling better handling of long-range dependencies with non-sequential structure.

\subsubsection{External Memory Augmentation}

LongMem introduced external memory modules through a decoupled architecture that kept the original language model frozen while adding a side network for memory retrieval and integration \citep{wang2023augmenting}. This design preserved the pre-trained model's capabilities while augmenting it with long-term memory spanning up to 65K tokens. The frozen backbone model served as a memory encoder, processing text to create memory representations, while an adaptive residual side network learned to retrieve and integrate relevant memories for the current context.

The decoupled design offered several advantages. By freezing the backbone model, LongMem avoided catastrophic forgetting that can occur when fine-tuning pre-trained models on long-context data. The side network could be trained specifically for the memory task without interfering with the model's existing capabilities. The external memory enabled caching of many-shot demonstrations for few-shot learning, allowing the model to access extensive in-context examples without exceeding its base context window.

Self-Extend \citep{jin2024llm} proposed a related but distinct approach that required no training whatsoever. Self-Extend employed bi-level attention where grouped attention with floor-based position remapping captured dependencies between distant tokens, while neighbor attention maintained fine-grained local relationships. This simple modification, requiring only four lines of code changes, achieved performance matching or exceeding fine-tuning methods on various long-context benchmarks. The success of Self-Extend demonstrated that context extension can sometimes be achieved through clever inference-time modifications to position encoding without the need for extensive retraining.

\subsection{Continual Pre-training for Context Extension}

While architectural and algorithmic innovations provide the mechanisms for processing long contexts, effective deployment of long-context capabilities in practice has increasingly relied on continual pre-training strategies. Most large language models undergo initial pre-training on large-scale corpora with relatively modest context windows, typically 2-4K tokens, due to the computational expense of training with longer sequences. Following this base pre-training, a separate context extension phase through continual pre-training on long-context data has emerged as standard practice for developing production long-context models.

\subsubsection{Data Construction and Training Strategies}

Research into optimal continual pre-training strategies has revealed that data composition critically influences long-context capabilities. Code repositories provide an excellent source of naturally long, highly structured data where dependencies span many lines. Books offer long-form narrative coherence that exercises different aspects of long-range understanding than code. Perhaps counterintuitively, including high-quality short data in the continual pre-training mixture proves essential for maintaining short-context performance, preventing the degradation that occurs when models are exclusively trained on long sequences.

Evidence suggests that training with sequence lengths exceeding the target evaluation length improves long-context performance, possibly because longer training sequences provide more opportunities to learn long-range dependencies and develop robust positional understanding. Practical training configurations vary by model scale. For 7B and 13B parameter models, training with 32,768-token sequences has proven effective, while larger 34B and 70B parameter models often use 16,384-token sequences due to memory constraints. Typical continual pre-training regimes process 400B additional tokens at extended sequence lengths, representing a substantial but manageable fraction of the original pre-training data volume.

Recent open-source frontier models including Llama-3.1 and Jamba have employed continual pre-training strategies that combine long-context continued training with subsequent supervised fine-tuning on instruction-following data. An interesting finding from these efforts indicates that using predominantly short instruction datasets during the supervised fine-tuning phase can yield strong performance on long-context tasks, provided the model has received adequate long-context continual pre-training. This observation suggests that long-context understanding and instruction-following capabilities can be developed somewhat independently, with continual pre-training establishing the long-context capacity and supervised fine-tuning adapting it to instruction-following without requiring extensive long-form instruction data.

Code Llama exemplifies the continual pre-training approach, having been developed through continued training of LLaMA 2 checkpoints on an additional 400B tokens formed into long training sequences. The domain specificity of code, with its structured syntax and long-range dependencies between definitions and uses, makes it particularly amenable to long-context modeling. LongLLaMA-Code, a variant of Code Llama, incorporated Focused Transformer (FoT) methods during fine-tuning to further extend context handling, demonstrating that domain-specific continual pre-training can be combined with architectural modifications for enhanced long-context capability.

\subsubsection{LongAlign and Comprehensive Alignment Recipes}

LongAlign established the first complete recipe for long-context alignment, addressing both the continual pre-training phase and the subsequent supervised fine-tuning required to align long-context models with user intentions \citep{longalign2024}. The continual pre-training component involved expanding the RoPE base frequency by 200$\times$ (from 10K to 2M), conducting continual training on sequences up to 64K tokens for a total of 10B tokens, and developing loss-weighting strategies for packed training with varied lengths. These techniques enabled efficient utilization of training data spanning diverse lengths while maintaining balanced learning across different sequence lengths.

LongAlign's key contribution included the LongAlign-10k dataset, containing 10,000 long instruction examples ranging from 8K to 64K tokens in length. This dataset was constructed using collected long sequences from nine diverse sources and applying the Self-Instruct approach with a long-context LLM (Claude 2.1) to generate instruction-following examples. The availability of this dataset addressed a critical gap in long-context alignment, as prior work lacked sufficient long-form instruction data for effective supervised fine-tuning.

Training efficiency improvements constituted another important aspect of LongAlign. Packing with loss weighting enabled processing multiple sequences of varying lengths within a single training batch, with loss weights ensuring that shorter sequences did not dominate the training signal. Sorted batching grouped similar-length sequences together to minimize padding waste and improve computational efficiency. These techniques significantly reduced the computational cost of long-context fine-tuning while improving final model quality.

LongAlign introduced LongBench-Chat as an evaluation benchmark specifically designed for assessing instruction-following capability on queries requiring 10K-100K token contexts. This benchmark filled a gap in evaluation methodology, as prior benchmarks focused primarily on retrieval tasks or short-context instruction following. The comprehensive nature of LongAlign, spanning continual pre-training, dataset construction, training efficiency, and evaluation methodology, established it as a foundational reference for practitioners developing long-context models. The work demonstrated conclusively that effective long-context capability requires coordinated attention to both continual pre-training for context extension and proper alignment through supervised fine-tuning with long instruction data.

\subsubsection{Efficient Extension Methods}

Recent work has focused on reducing the computational cost of continual pre-training for context extension. E²-LLM (Efficient and Extreme Length Extension) achieved dramatic reductions in training costs through a dual RoPE augmentation strategy that trained exclusively on short sequences while still enabling extreme length extension \citep{xnhyacinth2024e2llm}. By scaling and adjusting position indices across samples during training on 4K sequences, E²-LLM eliminated the need to collect and train on long-context data entirely. This approach reduced continual pre-training costs by orders of magnitude while supporting variable context windows at evaluation time, demonstrating that careful manipulation of position encodings during training can substitute for actual long-sequence training in some scenarios.

CLEX (Continuous Length Extrapolation) modeled context extension dynamics through ordinary differential equations over length scaling factors \citep{chen2024clex}. This formulation enabled smooth extrapolation to 4-8$\times$ the training context length without performance degradation, with a 4K-trained model achieving competitive performance at 32K on the LongBench benchmark. CLEX required minimal additional training or inference latency and integrated seamlessly into RoPE-equipped models including LLaMA and GPT-NeoX. The ODE-based perspective provided theoretical grounding for understanding length extrapolation as a continuous process rather than a discrete scaling operation.

Data engineering for context extension has received substantial attention as researchers recognized that data quality and diversity often matter more than sheer quantity. Comprehensive studies of data engineering for 128K context scaling have examined training data construction, curriculum design, and evaluation methodology. These studies demonstrated that strategic data selection, incorporating diverse long-context sources and maintaining data quality standards, outperforms naive approaches that simply concatenate documents to create long training sequences.

\subsubsection{Training Curriculum and Scheduling}

SkyLadder introduced context window scheduling as a training-time optimization that progressively increases sequence length across training phases \citep{li2025skyladder}. Rather than training with a fixed maximum sequence length throughout the entire training run, SkyLadder begins with shorter sequences and gradually increases the context window as training progresses. This curriculum-based approach reduces total compute requirements because earlier training phases with shorter sequences are substantially cheaper, while still enabling the model to develop long-context capabilities through exposure to progressively longer sequences in later phases. Empirical results demonstrated that SkyLadder reduced pretraining compute while simultaneously improving final model quality, indicating that the curriculum learning provided better optimization trajectories than static long-context training.

Recent research has also examined the interaction between context length during supervised fine-tuning and resulting model behavior. Counterintuitively, supervised fine-tuning at extended context lengths was found to improve short-context performance in certain scenarios, contrary to the expectation that long-context fine-tuning would exclusively benefit long-context tasks. This bidirectional relationship between short and long-context capabilities during training suggests complex interactions in how models represent and utilize positional information, informing the design of more efficient multi-scale context training procedures.

Synthetic data generation for continued pre-training has emerged as a practical approach for resource-constrained settings where collecting natural long-context corpora proves difficult. Synthetic continued pre-training explores generation of training data specifically designed for length extension, reducing dependency on naturally occurring long sequences. While synthetic data cannot fully replace diverse natural data for general long-context capability, targeted synthetic generation can provide cost-effective training data for specific long-context tasks or domains where natural data is scarce.

\subsection{Emerging Techniques and Current Frontiers}

Beyond the established categories of context extension methods, emerging research explores novel approaches that combine insights from multiple paradigms or introduce entirely new perspectives on the problem.

\subsubsection{Adaptive and Learned Approaches}

Core Context Aware Attention (CCA-Attention) introduced globality-pooling attention where tokens communicate through a reduced-size set of core tokens, reducing computational complexity while preserving semantic understanding \citep{zhang2024core}. This approach balances efficiency gains with performance preservation by maintaining dense attention within local regions while using pooled global tokens for long-range communication. Similar ideas appear in tokenwise attention scaling, which proposes dynamic scaling of attention mechanisms on a per-token basis to maintain stability across varying context lengths, and context-aware routing mechanisms that use learned routing to distribute attention computation intelligently across the context.

TokenSelect applied learned scoring to identify important tokens for retention in the KV cache during inference, enabling efficient processing of ultra-long sequences without approximating full attention \citep{jiang2025tokenselect}. By selectively retaining only tokens deemed important by a learned scoring function, TokenSelect reduced memory requirements and inference latency while maintaining high task performance. This selective retention approach proved particularly effective for retrieval tasks where relevant information concentrates in small portions of the full context.

Extending LLM context with adaptive grouped attention proposed grouping tokens for attention computation with group sizes that adapt based on position and importance \citep{chang2025extending}. This enabled smooth context scaling with minimal training overhead through learned grouping strategies that automatically balanced computational efficiency and representational capacity. The adaptive nature allowed the model to allocate more attention computation to important regions while processing less critical regions more efficiently.

\subsubsection{Training Efficiency and Sample Efficiency}

Research has demonstrated that context extension can be achieved with remarkably limited training data when data is carefully selected. Work showing context window extension with merely 100 training samples illustrated that sample-efficient training is possible through optimal selection of training instances and careful optimization strategies. This dramatic reduction in data requirements enables rapid prototyping and experimentation with context extension techniques without the computational resources required for large-scale continual pre-training.

Partial context training explored the counterintuitive finding that efficient long-context training can be achieved through partial context masking during training, where the model is exposed to incomplete contexts rather than always receiving full context. This approach reduced memory requirements during training while paradoxically improving generalization to full contexts at inference time, possibly because partial contexts forced the model to develop more robust representations that could handle missing information.

Segmented base adjustment for RoPE proposed segment-wise adjustment of RoPE base frequencies rather than global scaling, enabling more flexible adaptation to different context regions \citep{kim2024extending}. By dividing input sequences into segments and applying optimized base adjustments per segment, this approach improved context extension efficiency and performance compared to uniform scaling approaches.

\subsubsection{Theoretical Foundations and Analysis}

Recent theoretical work has explored probabilistic frameworks for understanding context extension. Bayesian Attention proposed treating position uncertainty explicitly through Bayesian inference, providing a principled framework for length extrapolation \citep{chen2025bayesian}. This probabilistic perspective offered theoretical grounding for understanding when and why certain extrapolation strategies succeed, potentially guiding the development of more robust extension methods.

Meta-analyses examining the fundamental requirements for context scaling have argued that extending context requires rethinking attention mechanisms themselves, not merely adjusting position encodings. These analyses advocate for integrated approaches that combine position scaling, attention approximation, and memory strategies, recognizing that optimal long-context performance emerges from coordinated advances across multiple components of the architecture rather than optimization of any single component in isolation.

Work on sparse attention tradeoffs has provided decision frameworks for choosing between sparse and dense attention based on task requirements, sequence length, and hardware constraints. Comprehensive surveys of sparse attention patterns, including fixed patterns, learned patterns, and dynamic routing approaches, have organized the design space and provided guidance on when different approximation strategies prove effective.

\subsection{Summary and Current State}

Context window extension techniques have evolved rapidly from simple position interpolation approaches to sophisticated integrated systems combining frequency-aware scaling, memory-efficient computation, distributed processing, and targeted continual pre-training. The field has progressed from early 2K-4K token limitations to contemporary production deployments at 128K tokens and research demonstrations exceeding 2 million tokens. This thousand-fold increase in practical context capacity has been achieved through complementary advances across multiple fronts.

RoPE-based scaling methods, progressing from Position Interpolation through NTK-Aware scaling and YaRN to LongRoPE, have established that careful frequency-domain manipulation enables dramatic context extension with minimal fine-tuning. The key insight underlying these methods is that different frequency components of positional embeddings serve different representational purposes and must be scaled accordingly, with high frequencies preserving local coherence and low frequencies maintaining global structure. Alternative position encoding approaches like ALiBi have demonstrated that strong length extrapolation is possible without traditional positional embeddings, though with different tradeoffs regarding recency bias and task generality.

Architectural and system-level innovations, particularly FlashAttention and Ring Attention, have made long-context processing computationally feasible by addressing memory bandwidth bottlenecks and enabling distributed computation. These advances demonstrate that the quadratic complexity of attention, while fundamental to the algorithm, need not prevent processing of very long sequences when memory hierarchy and device communication are carefully optimized. The progression from FlashAttention through FlashAttention-2 and FlashAttention-3, each achieving substantial additional speedups through deeper hardware-software co-design, illustrates the continuing potential for system-level optimization.

Memory-augmented approaches and streaming methods have shown that explicit memory mechanisms, attention sinks, and block-level retrieval can complement standard attention for handling ultra-long contexts. The discovery of attention sinks in pre-trained models highlights the importance of understanding emergent behaviors in trained systems, while learned memory mechanisms like those in RMT demonstrate that combining transformer attention with recurrent state propagation can achieve linear scaling in computation.

Continual pre-training has emerged as essential practice for deploying long-context capabilities, with research establishing optimal data compositions, training curricula, and alignment strategies. The recognition that long-context capability requires both architectural support for processing long sequences and training procedures that expose models to long-range dependencies reflects a mature understanding of the multifaceted nature of context extension. Methods like LongAlign that provide comprehensive recipes spanning continual pre-training, dataset construction, and evaluation methodology have accelerated practical deployment by providing validated frameworks rather than requiring practitioners to discover effective strategies through trial and error.

Current challenges include maintaining short-context performance while extending long-context capability, a persistent issue addressed through progressive extension strategies, short-context recovery phases, and mixed-length training curricula. Computational constraints remain significant despite algorithmic advances, motivating continued work on sparse attention, efficient training methods, and inference-time optimization. Evaluation methodology has evolved to assess not merely retrieval capability but genuine long-context understanding, reasoning, and instruction following, with benchmarks like LongBench-Chat providing more comprehensive assessment of long-context competence.

The convergence of these techniques in modern production systems, which typically combine RoPE-based scaling with FlashAttention optimization and continual pre-training on long-context data, demonstrates that effective context extension requires coordinated advances across position encoding, attention computation, system optimization, and training methodology. No single technique proves sufficient in isolation; rather, the combination of complementary approaches yields practical long-context language models. As research continues to push toward million-token contexts and beyond, the integration of insights from theoretical analysis, empirical optimization, and system design promises further advances in humanity's ability to equip artificial intelligence systems with extended memory and reasoning capabilities over long sequences.
